\documentclass[10pt,a4paper]{article}
\usepackage[utf8]{inputenc}
\usepackage[T1]{fontenc}
\usepackage[spanish,es-nodecimaldot]{babel}
\usepackage{lmodern}
\usepackage[margin=1.85cm,headheight=14pt]{geometry}
\usepackage{xcolor}
\usepackage{booktabs,tabularx,array,longtable}
\usepackage{enumitem}
\usepackage{fancyhdr}
\usepackage{hyperref}
\usepackage{listings}
\usepackage{microtype}
\definecolor{navy}{HTML}{173F67}
\definecolor{blue}{HTML}{276FA8}
\definecolor{soft}{HTML}{F3F6F9}
\definecolor{line}{HTML}{D5DFE8}
\hypersetup{colorlinks=true,linkcolor=navy,urlcolor=blue}
\pagestyle{fancy}
\fancyhf{}
\lhead{\small Seminario de Programación - Periodo 2}
\rhead{\small Instituto Jorge Robledo}
\cfoot{\small \thepage}
\setlength{\parindent}{0pt}
\setlength{\parskip}{5pt}
\setlist[itemize]{leftmargin=5mm,itemsep=2pt,topsep=3pt}
\setlist[enumerate]{leftmargin=6mm,itemsep=3pt,topsep=3pt}
\newcommand{\ruleline}{\par\noindent\color{line}\rule{\linewidth}{0.6pt}\color{black}\par}
\newcommand{\statusbox}[2]{\noindent\fcolorbox{line}{soft}{\parbox{0.94\linewidth}{\textbf{#1}\\#2}}}
\newcommand{\scorebox}[2]{\fcolorbox{line}{soft}{\parbox{0.43\linewidth}{\centering\textcolor{navy}{\Large\textbf{#1}}\\[-1mm]\small #2}}}
\lstset{basicstyle=\ttfamily\small,backgroundcolor=\color{soft},frame=single,rulecolor=\color{line},breaklines=true,columns=fullflexible,showstringspaces=false}
\begin{document}
\begin{center}
{\color{navy}\Large\bfseries Proyecto de referencia desde cero}\\[2mm]
{\small Carrito web completo con HTML, CSS, JavaScript, DOM y localStorage}
\end{center}
\ruleline

\section*{Propósito}
Construir una aplicación que permita seleccionar productos, validar cantidades, calcular subtotales y total, guardar el pedido y mostrar un recibo. Este estándar se adapta a otras temáticas; no obliga a copiar el carrito.
\section*{Estructura mínima}
\begin{lstlisting}
mi-proyecto/
  index.html
  recibo.html
  css/styles.css
  js/products.js
  js/cart.js
  js/app.js
  js/receipt.js
  README.md
\end{lstlisting}
\section*{Criterios}
\begin{tabularx}{\linewidth}{>{\raggedright\arraybackslash}X X}
\toprule
\textbf{Elemento} & \textbf{Evidencia esperada}\\\midrule
HTML & header, main, section, formularios, labels y salida visible\\
JavaScript & objetos, arreglos, funciones, condiciones y cálculos\\
DOM y eventos & addEventListener, renderizado y mensajes\\
CSS & variables, Grid/Flexbox, responsive y estados\\
Persistencia & JSON.stringify/JSON.parse con localStorage\\
Validación & entradas válidas y errores visibles\\
Integración & entrada -> proceso -> estado -> salida\\
Git & README y commits pequeños\\
\bottomrule
\end{tabularx}

\newpage
\section*{1. Modelo de datos}
\begin{lstlisting}
export const products = [
  { id: "hamburguesa", name: "Hamburguesa", price: 15000 },
  { id: "pizza", name: "Pizza personal", price: 18000 },
  { id: "gaseosa", name: "Gaseosa", price: 5000 }
];
\end{lstlisting}
\section*{2. HTML principal}
\begin{lstlisting}
<header><h1>Carrito de restaurante</h1></header>
<main>
  <section><h2>Productos</h2><div id="product-list"></div></section>
  <section>
    <h2>Pedido</h2>
    <div id="cart-list"></div>
    <p>Total: <strong id="cart-total">$0</strong></p>
    <p id="form-message" role="status"></p>
    <button id="checkout-button" type="button">Finalizar</button>
  </section>
</main>
\end{lstlisting}
\section*{3. Funciones}
\begin{lstlisting}
const cart = [];
function addToCart(productId, quantity) {
  const amount = Number(quantity);
  if (!Number.isInteger(amount) || amount <= 0) {
    throw new Error("La cantidad debe ser un entero positivo");
  }
  const product = products.find(item => item.id === productId);
  if (!product) throw new Error("Producto no encontrado");
  const existing = cart.find(item => item.id === productId);
  if (existing) existing.quantity += amount;
  else cart.push({ ...product, quantity: amount });
}
function calculateTotal() {
  return cart.reduce((sum, item) => sum + item.price * item.quantity, 0);
}
\end{lstlisting}

\newpage
\section*{4. DOM y eventos}
\begin{lstlisting}
function renderCart() {
  cartList.innerHTML = "";
  for (const item of cart) {
    const row = document.createElement("article");
    row.textContent = `${item.name} - ${item.quantity * item.price}`;
    cartList.append(row);
  }
  totalOutput.textContent = `$${calculateTotal()}`;
}
productList.addEventListener("click", event => {
  const button = event.target.closest("[data-add-product]");
  if (!button) return;
  try { addToCart(button.dataset.addProduct, 1); renderCart(); }
  catch (error) { message.textContent = error.message; }
});
\end{lstlisting}
\section*{5. Persistencia y recibo}
\begin{lstlisting}
localStorage.setItem("restaurant-cart", JSON.stringify(cart));
window.location.href = "recibo.html";

const saved = JSON.parse(localStorage.getItem("restaurant-cart") || "[]");
\end{lstlisting}
\section*{6. CSS}
\begin{lstlisting}
:root { --ink:#172033; --surface:#fff; --canvas:#f3f6f9;
  --primary:#1f5f96; --line:#d6e0e8; --radius:12px; }
.product-grid { display:grid;
  grid-template-columns:repeat(auto-fit,minmax(210px,1fr)); gap:1rem; }
button:focus-visible { outline:3px solid #78aeda; }
button:disabled { opacity:.55; cursor:not-allowed; }
@media (max-width:650px) { main { padding:.75rem; } }
\end{lstlisting}

\newpage
\section*{7. README y commits}
\begin{lstlisting}
# Nombre del proyecto
## Integrantes y roles
## Objetivo
## Como ejecutar
## Flujo principal
## Validaciones
## Estado actual y pendientes
\end{lstlisting}
\begin{lstlisting}
feat: renderiza catalogo de productos
feat: agrega productos y calcula total
fix: valida cantidades enteras positivas
feat: guarda datos y genera recibo
style: crea componentes y responsive
docs: documenta integrantes y ejecucion
\end{lstlisting}
\section*{8. Pruebas}
\begin{enumerate}
\item Abrir sin errores 404.
\item Probar entrada válida y errores.
\item Confirmar cálculos y actualización del DOM.
\item Recargar y verificar persistencia.
\item Revisar móvil, teclado, foco visible y mensajes.
\item Confirmar integrantes, roles y autoría.
\end{enumerate}
\section*{Meta del referente}
Jerónimo Rodríguez Peña debe cerrar persistencia y recibo, añadir validaciones explícitas, responsive, variables CSS y documentación. Ser referente significa avanzar hacia un flujo completo y explicable, no conservar el mismo puntaje.
\end{document}
